\documentclass[12pt]{article}
\usepackage[utf8]{inputenc}
\usepackage[margin=0.75in]{geometry}
\usepackage{amsmath}
\usepackage{amssymb}
\usepackage{amsthm}
\usepackage{graphicx} % Required for inserting images
\usepackage{array}
\usepackage{tikz-cd}


\title{MAT1100 Homework 2}
\author{Aaron Avram}
\date{October 6 2026}

\begin{document}
\newtheorem{definition}{Definition}
\newtheorem{theorem}{Theorem}
\newtheorem{claim}{Claim}
\newtheorem{corollary}{Corollary}

\maketitle

\begin{align*}
    \textbf{Question 3}
\end{align*}

Arithmetic for the permutations is done mod p unless explicitly stated otherwise.
\begin{enumerate}
    \item[(a)]
    \begin{proof}
        Let $p$ be prime and let $\mathbb F_p = \mathbb Z/p\mathbb Z$ which is a field.
        Consider $\tau(x) = x+1$ which is an element of $S_{\mathbb F_p}$ and let $P = \langle \tau \rangle$.
        Suppose we define $N = \{ \sigma \in S_{\mathbb F_p}: \sigma P \sigma^{-1} = P \}$.
        To find all $\sigma \in N$ it suffices to find all $\sigma \in S_{\mathbb F_p}$ such that $\sigma \tau \sigma^{-1} = \tau^a$
        as $P$ is generated by $\tau$. Note that $\tau$ has cycle representation $(0,1,\ldots, p-1) = \Big(0, \tau(0), \tau^2(0), \ldots, \tau^{p-1}(0)\Big)$
        and so
        \[ \sigma \tau \sigma^{-1} = \Big(\sigma(0),\sigma(1),\ldots,\sigma(p-1)\Big) \]
        now $\sigma(0)$ could be any element of $\mathbb F_p$. However it is clear that
        \[ \tau^a(x) = x + a \]
        and hence $\tau^a$ has cycle decomposition $(0, ka, 2ka, \ldots (p-1)a) = \Big(0, \tau^a(0), \tau^{2a}(0), \ldots, \tau^{p-1}(0)\Big)$
        so if
        \[ \sigma \tau \sigma^{-1} = \tau^a \]
        we have
        \[ \Big(\sigma(0),\sigma(1),\ldots,\sigma(p-1)\Big) = \Big(0, \tau^a(0), \tau^{2a}(0), \ldots, \tau^{p-1}(0)\Big) \]
        and hence if $\sigma(0) = ka$ then $\sigma(1) = (k+1)a$.
        $\sigma$ is now completely determined as this implies $\sigma(i) = (i + k)a$.
        Therefore all such possible $\sigma$ are determined by whether they can map $0$ and the possible values of $a$.
        $\sigma(0) \in \{ 0, \ldots, p-1 \}$ and $a \in \{ 1, \ldots, p-1 \}$ where $a \neq 0$ as no permutation
        is conjugate to the identity.

        Now this implies that $|N| = p(p-1)$ with explicit representation
        \[ N = \{ \sigma \in S_{\mathbb F_p}: \exists a \in \{ 1, \ldots, p-1 \} \text{ s.t. } \forall i \in \{ 0, \ldots, p-1\}, \sigma(i+1) - \sigma(i) = a \} \]
        i.e.
        \[ N = \{ \sigma(x) = ax + b \in S_{\mathbb F_p}: a \in \mathbb F_p^\times, b \in \mathbb F_p \} \]
        to show $N$ is a subgroup of $S_{\mathbb F_p}$ it suffices to show that it is non empty
        and closed under the group operation as $S_{\mathbb F_p}$ is finite.
        Clearly $e \in N$ so $N \neq \varnothing$. Now if $\sigma, \rho \in N$ then
        \[ (\sigma \rho)P(\sigma \rho)^{-1} = \sigma(\rho P \rho^{-1})\sigma^{-1} = \sigma P \sigma^{-1} = P \]
        and thus $\sigma\rho \in N$. Hence $N \leq G$.

        Next we show that $P \lhd N$. Clearly $P \leq N$ as conjugating $P$
        by itself is an automorphism of $P$. Normality is clear as if $\sigma \in N$
        then $\sigma P \sigma^{-1} = P$ by definition.

        Now to characterize $N$ further, note that the characterization
        \[ N = \{ \sigma \in S_{\mathbb F_p}: \exists a \in \{ 1, \ldots, p-1 \} \text{ s.t. } \forall i \in \{ 0, \ldots, p-1\}, \sigma(i+1) - \sigma(i) = a \} \]
        gives a bijection of sets
        from $N$ to $\mathbb F_p \times \mathbb F_p^\times$ given by the map
        \[ f: \sigma_{b,a} \to (b, a) \]
        where $\sigma_{b, a}$ is the permutation given by $\sigma(x) = ax + b$.

        Now I claim that $N \cong \mathbb F_p \rtimes_\varphi \mathbb F_p^\times$ with
        multiplication in the semidirect product defined as
        \[ (b, a) \cdot (b', a') = (b + ab', aa') \]
        where $\varphi: \mathbb F_p^\times \to \operatorname{Aut}(\mathbb F_p)$ is
        the identity as the automorphisms of $\mathbb F_p$ are precisely multiplication
        by a non-zero element (if $\phi$ is an automorphism it corresponds to multiplication by $\phi(1)$).

        Now we need only show that the $f$ above is a homorphism as it is already clearly a bijection.
        First note that
        \[ \sigma_{b, a} \circ \sigma_{b', a'}(x) = \sigma_{b, a}(a'x + b') = (aa'x + ab' + b) = \sigma_{ab' + b, aa'}(x) \]
        Therefore
        \[ f(\sigma_{b, a}\circ\sigma_{b', a'}) = f(\sigma_{ab' + b, aa'}) = (ab' + b, aa') = (a, b)(a', b') = f(\sigma_{a, b})f(\sigma_{a', b'}) \]
        as required. Now clearly we can identify $P$ with the functions $x + b$ for $b \in \mathbb F_p$
        therefore $P \cong \mathbb F_p$. Hence
        \[ N/P \cong \mathbb F_p^\times \]
    \end{proof}

    \item[(b)]
    \begin{proof}
        $|S_{\mathbb F_p}| = p!$ and so $p \mid |S_{\mathbb F_p}|$ but
        $p^2 \not\mid |S_{\mathbb F_p}|$ as otherwise this would imply $p \mid (p-1)!$
        and since $p$ is prime this means there exists $k < p$ such that $p \mid k$ which is clearly a contradiction.

        Now $n_p = 1 \mod p$ and furthermore $P$ from $(a)$ is clearly a Sylow p-subgroup
        of $S_{\mathbb F_p}$ with normalizer $N$.
        Therefore
        \[ n_p = [S_{\mathbb F_p} : P] = (p-2)! \]
        hence $S_{\mathbb F_p}$ has $(p-2)!$ Sylow p-subgroups.
        Therefore
        \[ 1 = (p-2)! \mod p \]
        and multiplying by $p-1$
        \[ p-1 = (p-1)! \mod p \]
        and hence
        \[ -1 = (p-1)! \mod p \]
    \end{proof}

    \item[(c)]
    \begin{proof}
        See (b).
    \end{proof}
\end{enumerate}

\end{document}